\documentclass[10pt,a4paper]{amsart}
\usepackage[margin=19mm]{geometry}
\usepackage{amsmath,amssymb,mathtools}
\usepackage[scr=rsfs,cal=euler]{mathalfa}
\usepackage{xfrac}
\usepackage[unicode,hidelinks,bookmarks=false]{hyperref}
\newcommand{\ie}{i.e.\ }
\newcommand{\cf}{cf.\ }
\newcommand{\resp}{respectively\ }
\newcommand{\Subsection}{\subsection}
\providecommand{\nobreakdash}{\nobreak\hbox{-}}
\setlength{\emergencystretch}{2em}
\multlinegap0pt
\usepackage{amssymb}
\usepackage{tikz}
\newtheorem{prop}{Proposition}
\title{Coexact $1$-form spectral gaps of hyperbolic rational homology spheres}
\author{Francesco Lin, Michael Lipnowski}
\date{Source: arXiv:2603.02093v1 (CC BY 4.0)}
\begin{document}
\maketitle

\noindent\emph{Source and licence.} Coexact $1$-form spectral gaps of hyperbolic rational homology spheres. Version arXiv:2603.02093v1 (CC BY 4.0), distributed by arXiv under the Creative Commons Attribution 4.0 International licence, http://creativecommons.org/licenses/by/4.0/, as recorded in the arXiv licence metadata for this exact version and shown on the abstract page https://arxiv.org/abs/2603.02093v1. Full licence notice: \texttt{licenses/arxiv-2603.02093-CC-BY-4.0.txt}.\par\smallskip

\noindent\textit{Editorial scope.} Counted unit: the article's prop mainprop (source0.tex lines 288--293). The article's complete original proof of it is delivered with it (source0.tex lines 295--314, one \texttt{proof} environment of 20 lines). No mathematics is written, repaired or re-derived: the statement and the proof are the article's own lines, in the article's own notation, and no retained line is substituted or re-flowed.  Retained uncounted as the source-local context the proof needs: lines 168--170; lines 243--245; lines 277--279; lines 281--283.  Nothing was omitted from any retained span, so the package carries no omission record.  The retained lines cite 2 works, whose entries are transcribed verbatim from the article's own bibliography source in the e-print and delivered with short editorial labels recorded in the item's reference mapping.  No retained line contains a figure environment, an image-inclusion command or drawing code, so no image is delivered and the package declares no asset dependency.  Editorial formatting only, in addition: an AMS \texttt{amsart} preamble that restates the article's own declarations and macros used by the retained lines, namely the environments \texttt{prop} on separate counters, together with the standard packages those lines need.  The retained environments print in this document as Proposition 1 in order of appearance, whereas the article prints the same items with the numbering of its own full document.  The complete native source of the article as distributed in the arXiv e-print for this exact version is bundled unchanged as \texttt{sources/195597/source0.tex}.
\begin{equation}\label{conj}
I\circ\varphi\circ I^{-1}=\varphi^{-1}
\end{equation}

\begin{prop}\label{gap}
Suppose that no eigenvalue of $\varphi_*$ has absolute value $1$. Then $\{\lambda_1^*(Y_{\varphi,n})\}$ is uniformly bounded below.
\end{prop}

\begin{equation}\label{twistedhodge}
\Omega^1(L_{\psi})=d\Omega^0(L_{\psi})\oplus d^*\Omega^2(L_{\psi}),
\end{equation}

\begin{equation}\label{hodge}
\Omega^1=d\Omega^0 \oplus\mathbb{C}\oplus d^*\Omega^2
\end{equation}

\begin{prop}\label{mainprop}
In the setup of Proposition \ref{gap}, the quantity $\delta(\varphi)$ is an actual minimum which is \textbf{strictly positive}, and
\begin{equation*}
\mathrm{lim}_{n\rightarrow\infty} \lambda_1^*(Y_{\varphi,n})=\mathrm{inf}\{\lambda_1^*(Y_{\varphi,n})\}=\delta(\varphi)>0.
\end{equation*}
\end{prop}

\begin{proof} We study the symmetries of the cyclic covers $M_{\varphi,n}$ of $M_\varphi$. The condition (\ref{conj}) implies that $M_{\varphi,n}$ admits an isometric action of the dihedral group with $2n$ elements
\begin{equation*}
\mathsf{D}_n=\left\{r,s\mid r^n=1, s^2=1, srs=r^{-1}\right\},
\end{equation*}
where $s$ corresponds to the special involution $\iota$, and $r$ corresponds to the covering action by $\mathbb{Z}/n$. Considering the interaction of the latter with the eigenspace decomposition of coclosed $1$-forms, we obtain the identification (preserving the action of the twisted Hodge Laplacians)
\begin{equation}\label{cover}
d^*\Omega^2(M_{\varphi,n})=\bigoplus_{\zeta^n=1}d^*\Omega^2(L_\zeta)
\end{equation}
of coexact $1$-forms on $M_{\varphi,n}$ in terms of twisted coexact $1$-forms on the base manifold $M_\varphi$; let us analyze the latter first.
\par
The spectrum of $\Delta_\psi$ on $\Omega^1(L_\psi)$ varies continuously in the twisting $\psi\in U(1)$ \cite{Kat}. In fact, the direct sum decomposition (\ref{twistedhodge}) implies (for example via their min-max characterization) that the spectra of $\Delta_\psi$ on exact and coexact (twisted) $1$-forms are also both continuous in the twisting $\psi\in U(1)$ when $\psi\neq 1$. The key observation is that the $0$-eigenvalue of the untwisted Laplacian (corresponding to $b_1(M_\varphi) = 1$ as in (\ref{hodge})) will correspond under continuous deformation for twisting $\psi$ near $1\in U(1)$ to a very small eigenvalue on exact, rather than coexact, $1$-forms. This can be seen for example by noticing that the analogous decomposition to (\ref{cover}) holds for exact $1$-forms on $M_{\varphi,n}$, and for $n$ very large the cyclic cover $M_{\varphi,n}$ will have a very small spectral gap on functions (hence on exact $1$-forms) because it will have a very small Cheeger constant \cite{Cha}.
\par
This argument shows that the infimum of the spectral gap of $\Delta_{\psi}$ on coexact $1$-forms for $\psi$ in a small open neighborhood $O$ of $1\in U(1)$ is strictly positive. Finally, the infimum of this spectral gap for $\psi$ in the compact set $U(1)\setminus O$ is also strictly positive by continuity of the spectrum and the fact that $0$ does not belong in the spectrum of $\Delta_\psi$ by (\ref{twistedhodge}). In fact, continuity of the spectrum also implies that the infimum $\delta(\varphi)>0$ over $U(1)$ is realized.
\par
To conclude the proof, we consider the full action of the dihedral group $\mathsf{D}_n$ on $M_{\varphi,n}$. Again, $\mathsf{D}_n$ naturally acts on the eigenspaces; furthermore, the eigenforms that descend to the quotient $Y_{\varphi,n}=M_{\varphi,n}/\iota$ are exactly those fixed by $s$; this readily implies that
\begin{equation*}
\lambda_1^*(Y_{\varphi,n})\geq \lambda_1^*(M_{\varphi,n}) \text{ for all }n.
\end{equation*}
Suppose $\alpha$ is a coexact $1$-form with eigenvalue $\lambda$ on $M_{\varphi,n}$ on which $r$ acts as an $n$th root of unity $\zeta\neq \pm1$. This means that under the decomposition (\ref{cover}), $\alpha$ corresponds to a $\Delta_\zeta$-eigenform in the summand $d^*\Omega^2(L_\zeta)$. Then $s(\alpha)$ is also a coexact $1$-form with eigenvalue $\lambda$, and $r$ acts on it $\zeta^{-1}$. Therefore $s(\alpha)$ is different from $-\alpha$, and $s(\alpha)+\alpha$ is a non-zero $s$-invariant coexact $1$-form with eigenvalue $\lambda$ on $M_{\varphi,n}$, so $\lambda$ belongs to the coexact $1$-form spectrum of $Y_{\varphi,n}$. By the decomposition (\ref{cover}) and the continuity of the spectrum, the result follows.
\end{proof}

\begin{thebibliography}{99}
\bibitem[Cha]{Cha} Isaac Chavel. \newblock {\em Eigenvalues in {R}iemannian geometry}, volume 115 of {\em Pure and Applied Mathematics}. \newblock Academic Press, Inc., Orlando, FL, 1984. \newblock Including a chapter by Burton Randol, With an appendix by Jozef Dodziuk.
\bibitem[Kat]{Kat} Tosio Kato. \newblock {\em Perturbation theory for linear operators}. \newblock Classics in Mathematics. Springer-Verlag, Berlin, 1995. \newblock Reprint of the 1980 edition.
\end{thebibliography}
\end{document}
